\begin{chapintro}
検査の精度は、2×2 表から計算する感度と特異度で表します。陽性的中割合と陰性的中割合は有病割合によって変わるので、研究どうしを比べる指標には向きません。カットオフを動かすと感度と特異度は逆向きに動き、この性質が第 10 章以降の統計の考え方を決めます。分析性能と診断精度の違い、診断精度に関わる報告の指針も整理します。カットオフを動かして感度と特異度を計算する R のコードは発展の節に置いてあり、初読では飛ばせます。
\end{chapintro}

\section{2×2 表}
\label{sec:09-1}

診断精度研究では、対象者の全員に、精度を調べたい検査（対象検査）と、正解を決めるための検査（参照基準）の両方を行います。参照基準で標的疾患があると判定された人と、ないと判定された人に分け、それぞれで対象検査が陽性だった人と陰性だった人を数えると、4 つの升目を持つ表ができます。これが 2×2 表で、検査のメタ解析で研究ごとに取り出すデータはこの 4 つの数だけです。

第 2 部の例題は、発熱して受診した成人で、プロカルシトニン（PCT）が菌血症をどれだけ正しく見分けるかを調べた架空の研究 10 本です。参照基準は血液培養で、培養が陽性なら菌血症あり、陰性なら菌血症なしとします。最初の研究である架空の Kato 2010 は、救急外来の患者 355 人を対象にし、カットオフ 0.5 ng/mL 以上を陽性としました。この研究の 2×2 表を図 9.1 に示します。

\begin{kkcenter}
\begin{kkfitpic}
\begin{tikzpicture}[
  cell/.style={draw=kkaccent, line width=.5pt, minimum width=24mm, minimum height=13mm,
               align=flush center, font=\small},
  head/.style={minimum width=24mm, minimum height=8mm, align=flush center, font=\small\gtfamily}]
  \node[head] at (2.4,1.05) {菌血症あり};
  \node[head] at (4.8,1.05) {菌血症なし};
  \node[head, text width=20mm] at (-0.2,0) {PCT 陽性\\{\footnotesize 0.5 以上}};
  \node[head, text width=20mm] at (-0.2,-1.3) {PCT 陰性\\{\footnotesize 0.5 未満}};
  \node[cell, fill=kkpale] at (2.4,0) {真陽性 TP\\ 28};
  \node[cell] at (4.8,0) {偽陽性 FP\\ 61};
  \node[cell] at (2.4,-1.3) {偽陰性 FN\\ 17};
  \node[cell, fill=kkpale] at (4.8,-1.3) {真陰性 TN\\ 249};
  \node[font=\footnotesize] at (2.4,-2.3) {計 45 人};
  \node[font=\footnotesize] at (4.8,-2.3) {計 310 人};
  \node[font=\footnotesize, align=left, anchor=west] at (6.2,0) {計 89 人};
  \node[font=\footnotesize, align=left, anchor=west] at (6.2,-1.3) {計 266 人};
  \node[font=\scriptsize\gtfamily, text=kkaccent] at (3.6,1.75) {参照基準（血液培養）};
\end{tikzpicture}
\end{kkfitpic}
\kkcap{図 9.1　架空の Kato 2010 の 2×2 表（カットオフ 0.5 ng/mL、355 人）}
\end{kkcenter}

表の縦の列は参照基準による分け方で、横の行は対象検査の結果です。この並べ方は論文によって入れ替わることがあり、行と列を取り違えると感度と陽性的中割合を取り違えることになります。論文から数を取り出すときは、まず列と行が何で分けられているかを確かめます。

\section{感度と特異度}
\label{sec:09-2}

感度は、病気のある人のうち検査が陽性になった人の割合です。Kato 2010 では菌血症のある 45 人のうち 28 人が陽性だったので、感度は 28 ÷ 45 で 0.62 になります。特異度は、病気のない人のうち検査が陰性になった人の割合で、菌血症のない 310 人のうち 249 人が陰性だったので、249 ÷ 310 で 0.80 になります。

特異度の裏側にあたる量が偽陽性割合で、病気のない人のうち誤って陽性になった人の割合です。1 から特異度を引いた値と同じで、Kato 2010 では 61 ÷ 310 で 0.20 です。偽陽性割合は第 10 章以降の図の横軸に毎回登場するので、ここで名前を覚えておいてください。

感度と特異度は、どちらも分母が「参照基準で決まった病気の有無」の片側だけで閉じています。感度は病気のある人の列の中だけで、特異度は病気のない人の列の中だけで計算します。そのため、研究に病気のある人が多く含まれていても少なく含まれていても、検査の性質が同じなら感度と特異度は同じ値になると期待できます。研究どうしを比べたり統合したりする指標に感度と特異度を使うのは、この性質があるからです。ただし、この期待は「どの研究でも同じような患者に同じように検査をしている」という前提の上に立っており、現実には対象患者の重症度や測定法の違いで感度と特異度も変わります（Leeflang 2013）。その点は\kkref{sec:10-1}と\kkref{sec:14-1}で扱います。

\section{的中割合は有病割合で変わる}
\label{sec:09-3}

検査を受ける人の立場では、知りたいのは感度よりも「陽性と言われたとき、本当に病気である見込みはどのくらいか」です。これが陽性的中割合で、陽性になった人のうち本当に病気のある人の割合です（一般には陽性的中率とも呼ばれますが、分母に時間を含まないので、本書では割合と呼びます）。Kato 2010 では陽性 89 人のうち 28 人が菌血症だったので、0.31 になります。陰性的中割合は、陰性になった人のうち本当に病気のない人の割合で、266 人のうち 249 人なので 0.94 です。

的中割合は、検査の性質だけでなく、検査を受ける集団に病気の人がどれだけいるか（有病割合）で大きく変わります。Kato 2010 の感度 0.62 と特異度 0.80 をそのまま使い、有病割合だけを変えて 1000 人に検査をしたと考えると、表 9.1 のようになります。

\begin{kkcenter}
\kkfit{\begin{tabular}{rrrrrrr}
\toprule
有病割合 & 真陽性 & 偽陰性 & 偽陽性 & 真陰性 & 陽性的中 & 陰性的中 \\
\midrule
5\% & 31 & 19 & 187 & 763 & 0.14 & 0.98 \\
10\% & 62 & 38 & 177 & 723 & 0.26 & 0.95 \\
25\% & 156 & 94 & 148 & 602 & 0.51 & 0.86 \\
50\% & 311 & 189 & 98 & 402 & 0.76 & 0.68 \\
\bottomrule
\end{tabular}}
\kkcap{表 9.1　感度 0.62、特異度 0.80 の検査を 1000 人に行ったときの内訳（人数は四捨五入）}
\end{kkcenter}

有病割合が 5\% の集団では、陽性になった人の 7 人に 1 人しか菌血症ではありません。同じ検査を有病割合 50\% の集団で使えば、陽性の 4 人に 3 人が菌血症です。救急外来の発熱患者と ICU の重症患者では有病割合が違うので、同じ PCT の結果でも意味が変わります。研究ごとに的中割合を集めて平均すると、検査の性質と集団の性質が混ざった値になってしまいます。検査のメタ解析で的中割合を直接統合しないのはこのためで、統合するのは感度と特異度にし、的中割合は使う場面の有病割合を決めてから計算し直します（\kkref{sec:15-1}）。有病割合を連続的に動かしたときの的中割合の曲線は、第 13 章で得る要約点を使って、\kkref{sec:15-1-2}で扱います。

\section{尤度比と診断オッズ比}
\label{sec:09-4}

感度と特異度を組み合わせた指標に、尤度比があります。陽性尤度比は、病気のある人で陽性になる割合（感度）を、病気のない人で陽性になる割合（偽陽性割合）で割った値です。Kato 2010 では 0.622 ÷ 0.197 で 3.16 になり、病気のある人は、ない人の約 3 倍陽性になりやすいことを表します。陰性尤度比は偽陰性の割合（1 − 感度）を特異度で割った値で、0.378 ÷ 0.803 で 0.47 です。

尤度比は、検査前の見込みを検査後の見込みに換算するときに使います。オッズの考え方は\kkref{sec:05-1-4}で説明しました。検査前の見込み（事前確率）をオッズに直して尤度比を掛けると、検査後のオッズになります。事前確率 10\% はオッズで 1 対 9 です。陽性なら 3.16 を掛けて 3.16 対 9、確率に戻すと 0.26 になります。陰性なら 0.47 を掛けて確率に戻すと 0.05 です。表 9.1 の有病割合 10\% の行では、陽性的中割合が 0.26、陰性的中割合が 0.95（1 から引いた 0.05 が、陰性でも病気である見込み）で、尤度比から求めた値と一致します。尤度比を使った計算と 1000 人の表は、同じことを別の形で書いたものです。検査前の確率を横軸に、検査後の確率を縦軸にとった曲線は、発展の\kkref{sec:15-1-3}で扱います。

診断オッズ比は、陽性尤度比を陰性尤度比で割った値で、Kato 2010 では 6.72 になります。2×2 表の対角線の積の比（28 × 249 を 61 × 17 で割る）としても計算でき、検査の良さを 1 つの数で表せるので、以前のメタ解析ではよく統合されてきました。感度のオッズ（感度 ÷ (1 − 感度)）を、偽陽性割合のオッズで割った値としても書けます。

しかし、診断オッズ比が同じでも、感度と特異度の組み合わせは違いえます。検査 X の感度が 0.9、特異度が 0.5 のとき、診断オッズ比は (0.9 ÷ 0.1) ÷ (0.5 ÷ 0.5) で 9 です。検査 Y の感度が 0.5、特異度が 0.9 のときも、(0.5 ÷ 0.5) ÷ (0.1 ÷ 0.9) で 9 になります。X は陰性なら病気を除外しやすく、Y は陽性なら病気を確定しやすい検査です。除外に使いたいのか確定に使いたいのかで求める性質は違うので、1 つの数だけでは臨床の判断に使えません。本書では、診断オッズ比は補助的な指標として扱います。

\section{カットオフを動かすと何が起きるか}
\label{sec:09-5}

PCT は連続量で測る検査で、陽性と陰性を分けるのはカットオフです。菌血症のある人とない人の PCT の値の分布を描くと、図 9.2 のように 2 つの山が重なります。この図は図解のために仮想の分布から描いたもので、研究のデータではありません。PCT の対数が正規分布に従うと仮定し、菌血症のない人は中央値 0.2 ng/mL、ある人は中央値 1.5 ng/mL としました。対数については\kkref{sec:05-1-2}で説明しました。対数をとった値の標準偏差は、ない人で 1.0、ある人で 1.2 です。標準偏差 1.0 は、中央値のおよそ 2.7 倍から 2.7 分の 1 までの幅に、約 3 分の 2 の人が入る広さです。カットオフ 2 ng/mL の線は図解のために加えたもので、第 2 部の研究が使うカットオフは 0.25、0.5、1 ng/mL の 3 つです。

\kkfig[.8\linewidth]{x2_dist_cutoff.pdf}{図 9.2　菌血症のある人とない人の PCT の分布（仮想の分布。点線は左から 0.25、0.5、1、2 ng/mL のカットオフ）}

カットオフより右側に入った人が陽性になります。菌血症のある人の山で右側に入った部分が真陽性で、その割合が感度です。菌血症のない人の山で右側にはみ出した部分が偽陽性で、その割合が偽陽性割合です。4 つのカットオフについて、山の面積の割合を計算した結果を表 9.2 に示します（計算の R コードは、この節の終わりの発展に置きます）。

\begin{kkcenter}
\begin{tabular}{rrr}
\toprule
カットオフ（ng/mL） & 感度 & 特異度 \\
\midrule
0.25 & 0.93 & 0.59 \\
0.5 & 0.82 & 0.82 \\
1 & 0.63 & 0.95 \\
2 & 0.41 & 0.99 \\
\bottomrule
\end{tabular}
\kkcap{表 9.2　図 9.2 の仮想の分布での、カットオフごとの感度と特異度}
\end{kkcenter}

図の横軸は対数目盛なので、0.25、0.5、1、2 ng/mL は等間隔に並んでいます。菌血症のない人の山は 0.2 ng/mL 付近、ある人の山は 1.5 ng/mL 付近に頂点があり、2 つの山は 0.25 から 1 ng/mL のあたりで大きく重なっています。点線を右へ動かすと、菌血症のない人の山からはみ出す部分が減るので特異度は 0.59 から 0.99 まで上がります。同時に、菌血症のある人の山の右側に残る部分も減るので、感度は 0.93 から 0.41 まで下がります。左へ動かせば逆になります。2 つの山が重なっている限り、カットオフをどこに置いても、感度と特異度を同時に上げることはできません。

カットオフを決めるとは、2 つの山の重なりのどこに線を引くかを決めることです。基準範囲は健康な人の山だけを見て決めますが、カットオフは 2 つの山の両方を見て、見逃しと誤りのどちらを減らしたいかで決めます。感度と特異度が 0.82 で等しくなる 0.5 ng/mL は、この仮想の分布に置いた数字から偶然そうなっただけで、実際の検査で 0.5 ng/mL が釣り合う点だという意味ではありません。Kato 2010 のカットオフ 0.5 ng/mL での感度 0.62、特異度 0.80 と値が違うのも、図の分布が研究データではないためです。

カットオフを少しずつ動かしながら、横軸に偽陽性割合、縦軸に感度をとって点を打っていくと、図 9.3 の ROC 曲線ができます。先の 4 つのカットオフは、曲線の上の 4 つの点にあたります。

\kkfig[.65\linewidth]{x2_roc_single.pdf}{図 9.3　図 9.2 の分布から描いた ROC 曲線。点の横の数字はカットオフ（ng/mL）}

横軸は偽陽性割合（1 − 特異度）、縦軸は感度で、曲線上の 1 点が 1 つのカットオフに対応します。カットオフ 2 ng/mL の点は左下（偽陽性割合 0.01、感度 0.41）にあり、0.25 ng/mL の点は右上（偽陽性割合 0.41、感度 0.93）にあります。カットオフを下げるほど、曲線の上を右上へ進みます。点線の対角線は、検査の結果が菌血症と無関係なときの線で、曲線がこの線から左上へ離れているほど、2 つの山の重なりが小さいことを表します。

2 つの山の重なりが小さい検査ほど曲線は左上の角に近づき、重なりが完全なら曲線は対角線になります。曲線の下の面積（AUC）は、役に立つ検査なら 0.5 から 1 の値をとり、検査の識別能を 1 つの数で表します。ただし、診断オッズ比と同じく、AUC が同じでも曲線のどこを使うか（どのカットオフを使うか）で臨床の意味は違います。曲線は連続した線ですが、実際に使うのは 1 つのカットオフの 1 点です。

感度と特異度はカットオフを介して逆向きに動きます。メタ解析で集める研究は、それぞれ別のカットオフを使っていることが多く、研究ごとの感度と特異度は 1 本の ROC 曲線の上のあちこちに散らばります。研究どうしで感度が違って見えても、検査の性能が違うのではなく、カットオフが違うだけのことがあります。これを閾値効果と呼び、\kkref{sec:10-1}の出発点になります。

\subsection{\hatten 分布と ROC 曲線を R で描く}
\label{sec:09-5-1}

図 9.2 は R の dnorm() で描けます。dnorm(log(x), m, s) は、PCT が x のところの山の高さを返します。横軸は対数目盛にするので、plot() に log = "x" を指定します。

\begin{rcode}
# 対数をとった PCT が正規分布に従う仮想の集団
m1 <- log(1.5); s1 <- 1.2    # 菌血症の人
m0 <- log(0.2); s0 <- 1.0    # 菌血症でない人
x  <- exp(seq(log(0.01), log(50), length.out = 500))

plot(x, dnorm(log(x), m0, s0), type = "l", log = "x",
     xlab = "PCT（ng/mL、対数目盛）", ylab = "分布")
lines(x, dnorm(log(x), m1, s1), col = "red")
cuts <- c(0.25, 0.5, 1, 2)   # 4 つのカットオフ
abline(v = cuts, lty = 3)
\end{rcode}

pnorm() は、分布の左側の面積を返します。pnorm(log(k), m, s) は、PCT がカットオフ k 未満の人の割合です。

\begin{rcode}
# 菌血症の人でカットオフ以上の割合
sens <- 1 - pnorm(log(cuts), m1, s1)
# 菌血症でない人でカットオフ未満の割合
spec <- pnorm(log(cuts), m0, s0)
data.frame(cuts, sens = round(sens, 2), spec = round(spec, 2))
\end{rcode}

\begin{rout}
  cuts sens spec
1 0.25 0.93 0.59
2 0.50 0.82 0.82
3 1.00 0.63 0.95
4 2.00 0.41 0.99
\end{rout}

図 9.3 の ROC 曲線は、カットオフを細かく動かして同じ計算を繰り返し、偽陽性割合と感度を線で結んだものです。

\begin{rcode}
cc  <- exp(seq(log(0.005), log(200), length.out = 400))
fpr <- 1 - pnorm(log(cc), m0, s0)   # 偽陽性割合
tpr <- 1 - pnorm(log(cc), m1, s1)   # 感度
plot(fpr, tpr, type = "l", xlim = c(0, 1), ylim = c(0, 1),
     xlab = "1 − 特異度（偽陽性割合）", ylab = "感度")
abline(0, 1, lty = 3)               # 当て推量と同じ対角線
points(1 - spec, sens, pch = 19)    # 4 つのカットオフの点
text(1 - spec, sens, cuts, pos = 4)
\end{rcode}

\section{分析性能と診断精度、基準範囲とカットオフ}
\label{sec:09-6}

検査技師にとって検査の性能といえば、まず併行精度や室間再現精度、正確さ、直線性、検出限界といった分析性能です。分析性能は「同じ検体を測ったら同じ値が出るか、真の値に近い値が出るか」を表します。これに対して、感度と特異度が表すのは診断精度で、「その値で病気のある人とない人を分けられるか」です。分析性能は診断精度の前提であり、精度の悪い測定では診断精度も下がりますが、分析性能が良くても、病気のある人とない人の値がもともと重なっていれば診断精度は上がりません。

PCT のメタ解析では、分析性能の違いが研究間のばらつきの原因になりえます。測定試薬や機器が違えば同じ検体でも値が少しずれ、同じ 0.5 ng/mL というカットオフでも実質的に違う位置で線を引いていることになります。検出限界付近の低いカットオフ（0.25 ng/mL など）を使う研究では、低濃度域の測定のばらつきがそのまま判定のばらつきになります。検体の種類（血清か血漿か）、採血のタイミング（発熱から何時間後か）も値を動かします。PCT は感染の初期には上がりきっていないことがあるので、受診直後の 1 回の採血か、6 時間後の再検かで、同じ患者でも陽性と陰性が入れ替わります。論文から研究の特徴を取り出すときに、測定法と採血の時点を必ず項目に入れておくのはこのためです（\kkref{sec:12-1}）。

基準範囲とカットオフの違いにも注意が要ります。基準範囲は、健康な人の測定値の分布の中央 95\% を示すもので、病気を見分けるために決めた値ではありません。カットオフは、病気のある人とない人を分けるために、ある目的（見逃しを減らしたい、無駄な治療を減らしたい）に沿って選んだ値です。基準範囲の上限をそのままカットオフにした研究と、ROC 曲線を見て最も成績の良い点を選んだ研究とでは、カットオフの決め方の意味が違います。後者のように、集めたデータそのものを見てカットオフを選ぶと、その研究の中では精度が高めに出ます。これはバイアスのリスクの評価で問題になります（表 12.1、\kkref{sec:12-5}）。

\section{報告と評価のための指針}
\label{sec:09-7}

診断精度に関わる指針には、似た名前のものがいくつかあり、役割が違います。
\begin{description}
\item[STARD 2015] 個々の診断精度研究を報告するときの指針です（Bossuyt 2015）。研究を行う人が論文を書くときに使います。メタ解析の報告には使いません。
\item[PRISMA-DTA] 診断精度研究の系統的レビューとメタ解析を報告するときの指針です（McInnes 2018）。\kkref{sec:08-1}で説明した PRISMA 2020 の診断精度版にあたり、第 15 章で使います。
\item[QUADAS-2 と QUADAS-3] メタ解析に組み入れた個々の研究のバイアスのリスクと適用可能性を評価する道具です（Whiting 2011、Whiting 2026）。\kkref{sec:04-3}の RoB 2 にあたり、\kkref{sec:12-4}で使います。
\end{description}

STARD に沿って書かれた論文は、必要な情報が揃っていることが多く、2×2 表も作りやすくなります。とはいえ、STARD に準拠していることは研究の質が高いことを意味するわけではありません。報告が丁寧であることと、バイアスのリスクが低いことは別の問題で、後者は QUADAS で判定します。

\section{この章のまとめ}
\label{sec:09-8}

検査のメタ解析で研究ごとに集めるのは 2×2 表の 4 つの数で、統合するのは感度と特異度です。的中割合は有病割合で変わるので統合せず、使う場面の有病割合を決めてから計算します。尤度比と診断オッズ比は感度と特異度から計算できる補助の指標です。感度と特異度はカットオフを介して逆向きに動き、研究ごとにカットオフや測定法が違えば、精度の見かけの差が生じます。次の第 10 章では、この性質のために、介入のメタ解析の方法をそのまま当てはめられない理由を説明します。
